% Primary source: R. Chatterjee and P. Zhang, "On Ramsey indexes of graphs",
% J. Combin. Math. Combin. Comput. 129 (2026), 3--16, pp. 4--6.
% https://doi.org/10.61091/jcmcc129-01
% https://combinatorialpress.com/article/jcmcc/Volume%20129/on-ramsey-indexes-of-graphs.pdf
% Their [5] is G. Chartrand and P. Zhang, "The Ascending Ramsey Index of a
% Graph", Symmetry 15 (2023), 523, sec. 4; doi:10.3390/sym15020523.
% Open Problems Catalog: Q3768, "Full-length Ramsey chains in complete graphs".

\documentclass{article}
\usepackage{amsmath}
\begin{document}

\paragraph{Conjecture 2.6 (verbatim).}
For every integer $n \ge 3$, $AR(K_n) = n-1$.

\paragraph{Definitions from pp.\ 4--5, in notation.}
Let $G$ be a graph without isolated vertices, and let
$c:E(G)\to\{\mathrm{red},\mathrm{blue}\}$ be an edge colouring.
A Ramsey chain of length $\ell$ for $(G,c)$ consists of subgraphs
$G_1,\ldots,G_\ell$ satisfying
\begin{align*}
  |E(G_i)| &= i \quad (1\le i\le\ell),\\
  E(G_i)\cap E(G_j) &= \varnothing \quad (i\ne j),\\
  |c(E(G_i))| &= 1 \quad (1\le i\le\ell),\\
  G_i &\cong H_i\text{ for some subgraph }H_i\subseteq G_{i+1}
       \quad (1\le i<\ell).
\end{align*}
Here $H_i$ need not be induced, and the common colour within one link
need not be the common colour within another. The indexes are
\[
 AR_c(G)=\max\{\ell:(G,c)\text{ has a Ramsey chain of length }\ell\},
 \qquad AR(G)=\min_c AR_c(G).
\]

\paragraph{Counting consequence.}
For $G=K_n$ one has $\sum_{i=1}^{n-1}i=\binom n2$. Hence a chain of the
conjectured length partitions the edges.

\end{document}
